\documentclass[14pt,a4paper]{extreport}
% preamble.tex — спільна преамбула для всіх лабораторних робіт. Підключається одразу після \documentclass у кожному lab*/main.tex.
%
% ЗБІРКА: LuaLaTeX (не pdfLaTeX). Причина: пакет listings під pdfLaTeX не читає UTF-8 і зупиняє збірку на кириличних коментарях у лістингах програмного коду. LuaLaTeX обробляє UTF-8 нативно.
%
% ШРИФТ: Liberation Serif — метрично сумісний з Times New Roman і містить повний набір кириличних гліфів (TeX Gyre Termes кирилиці не має, через що український текст не відображався взагалі).

\usepackage{fontspec}
\setmainfont{Liberation Serif}
\setmonofont{DejaVu Sans Mono}[Scale=0.82]

\usepackage{polyglossia}
\setmainlanguage{ukrainian}
\setotherlanguage{english}

\usepackage{microtype}

% Поля, інтервал, відступи й заголовки нижче взяті з шаблону НУЛП, за яким написана бакалаврська робота (../bachelor-thesis/thesis/main.tex). Той самий формат приймає В.М. Шиманський і для магістерських, тож усі звіти тримаємо в ньому, а не в тому, що переказували на консультації.
\usepackage[left=25mm,right=15mm,top=15mm,bottom=23mm,footskip=1.5mm]{geometry}
\usepackage{setspace}
% Саме \setstretch{1.5}, а не \onehalfspacing: останній дає близько 1.25.
\setstretch{1.5}

\usepackage{graphicx}
\usepackage{float}
\usepackage{booktabs}
\usepackage{array}
\usepackage{tabularx}
\usepackage{longtable}
\usepackage{enumitem}
% labelindent вирівнює маркер списку з абзацним відступом — так у шаблоні НУЛП.
\setlist{nosep}
\setlist[itemize]{label=\textbullet,leftmargin=*,labelindent=1.27cm}
\setlist[enumerate]{leftmargin=*,labelindent=1.27cm}
\usepackage{indentfirst}
\usepackage{amsmath}

\usepackage{hyperref}
\hypersetup{colorlinks=true,linkcolor=black,urlcolor=blue,citecolor=black}

% xurl дозволяє \url розривати адресу в будь-якому місці, а не лише після скісної риски. Без нього довгий сегмент шляху (напр. «scopus---podrobnaya- instruktsiya-dlya-uchenogo») виходить за межі сторінки, бо дефіс для hyperref не є точкою розриву. Вантажиться після hyperref — перевизначає його \UrlBreaks.
\usepackage{xurl}

% ---------------------------- Список використаної літератури ----------------------------
% biblatex + biber зі стилем IEEE: посилання «[1]» у тексті, записи — за порядком першої згадки. Обрано IEEE, а не gost-numeric (ДСТУ 8302:2015), — так вирішив автор звітів після порівняння обох варіантів.
% Кожен звіт із джерелами сам підключає свій references.bib через \addbibresource і друкує список \printbibliography[heading=none] під власним \section. Звіт без джерел від завантаженого biblatex нічого не втрачає.
% Оточення списку — стандартне для стилю: мітка від лівого поля, текст опису вирівняний після неї. Не підганяємо його під абзацний відступ нумерованих переліків тексту (labelindent у \setlist вище) — список літератури є окремим розділом зі своїм стандартним оформленням.
\usepackage[backend=biber,style=ieee]{biblatex}
% biblatex за замовчуванням додає розтяжний проміжок після «/ . : ? =» в адресах, і з вимкненими переносами рядок вирівнюється розрідженням самої адреси («https : / / www . scopus»). Рвати адресу й так дозволяє xurl, тож проміжок лише шкодить.
\setlength{\biburlbigskip}{0mu}
% Без додаткового інтервалу між описами: розділ має йти суцільним текстом, як решта звіту. Так само в шаблоні бакалаврської.
\setlength{\bibitemsep}{0pt}
\setlength{\bibparsep}{0pt}

% ---------------------------- Лістинги програмного коду ---------------------------- Використано minted (через Pygments), а не listings: listings не вміє коректно обробляти багатобайтові символи й «висмикує» кириличні коментарі з коду, зліплюючи їх в один рядок. Потребує прапорця -shell-escape.
\usepackage{xcolor}
\usepackage[newfloat]{minted}
% labelsep=period дає «Рис. 2.1. Назва» замість «Рис. 2: Назва» — формат НУЛП.
\usepackage[labelsep=period]{caption}

% baselinestretch=1 обов'язковий: без нього лістинг успадковує півторашку основного тексту, і довгий фрагмент коду перестає вміщатися у сторінку. Усередині [H]-флоата це не помилка збірки, а нескінченний цикл порожніх сторінок — lab1 так намолотив 161 тисячу сторінок.
\setminted{
  baselinestretch=1,
  fontsize=\footnotesize,
  linenos=true,
  frame=single,
  framesep=4pt,
  breaklines=true,
  breakanywhere=true,
  tabsize=4,
  numbersep=7pt,
  bgcolor=gray!8,
}
\setminted[python]{python3=true}

\SetupFloatingEnvironment{listing}{name=Лістинг,listname=Перелік лістингів}

% ---------------------------- Нумерація та оформлення заголовків ----------------------------
\usepackage{titlesec}

\setcounter{secnumdepth}{3}
\renewcommand{\thesection}{\arabic{section}}
\renewcommand{\thesubsection}{\arabic{section}.\arabic{subsection}}
\renewcommand{\thesubsubsection}{\arabic{section}.\arabic{subsection}.\arabic{subsubsection}}

% Розділ по центру, 16 пт, великими літерами; слово «РОЗДІЛ» опущене, бо звіт не дисертація. Номер лишається: \numberwithin прив'язує до нього нумерацію рисунків («Рис. 2.1»), і без нумерованих розділів вона б зламалася. \MakeUppercase четвертим аргументом отримує текст заголовка, тому в джерелі заголовки пишуться звичайним регістром.
\titleformat{\section}[block]
  {\centering\bfseries\fontsize{16}{19.2}\selectfont}
  {\thesection.}{0.3em}{\MakeUppercase}
% Підрозділи — зліва з абзацного відступу, звичайним регістром.
\titleformat{\subsection}[block]{\bfseries}{\thesubsection.}{1em}{}
\titleformat{\subsubsection}[block]{\bfseries\itshape}{\thesubsubsection.}{1em}{}

% Незіркові \titlespacing — обов'язкові. Без них titlesec застосовує зірковий варіант, який прибирає абзацний відступ у першого абзацу після заголовка й перекриває indentfirst. Значення — типові для report.
\titlespacing{\section}{0pt}{20pt}{6pt}
\titlespacing{\subsection}{1.27cm}{6pt}{0pt}
\titlespacing{\subsubsection}{1.27cm}{6pt}{0pt}

\setlength{\parindent}{1.27cm}
\setlength{\parskip}{0pt}
\renewcommand{\arraystretch}{1.3}

% Переноси слів вимкнені — так у шаблоні НУЛП. Натомість абзац розганяється міжслівними пробілами, тому tolerance і emergencystretch підняті: без них рядки виходили б за праве поле.
\hyphenpenalty=10000
\exhyphenpenalty=10000
\tolerance=9999
\setlength{\emergencystretch}{3em}

% ---------------------------- Підписи до рисунків і таблиць ----------------------------
\captionsetup[figure]{name=Рис.,justification=centering,font={bf,it}}
\captionsetup[table]{name=Таблиця,justification=raggedleft,singlelinecheck=false,labelfont=bf,textfont=normalfont}

% Рисунки, таблиці й формули нумеруються в межах розділу: «Рис. 2.1», а не «Рис. 5». Так у шаблоні НУЛП, і так само виглядали приклади на консультації.
\numberwithin{equation}{section}
\numberwithin{figure}{section}
\numberwithin{table}{section}

% ---------------------------- Колонтитули ----------------------------
\usepackage{fancyhdr}
\pagestyle{fancy}
\fancyhf{}
\fancyfoot[R]{\thepage}
\renewcommand{\headrulewidth}{0pt}
\fancypagestyle{plain}{%
  \fancyhf{}
  \fancyfoot[R]{\thepage}
  \renewcommand{\headrulewidth}{0pt}
}

% titlepage.tex — титульний аркуш, спільний для всіх дисциплін.
% Потребує визначених у subject.tex команд:
%   \SubjectName, \LecturerName, \StudentName, \StudentGroup, \AcademicYear
%
% Використання у тексті: \MakeLabTitle{1}{Тема роботи}
%
% Рід роботи. Дисципліни з практичними роботами перевизначають його після
% % preamble.tex — спільна преамбула для всіх лабораторних робіт. Підключається одразу після \documentclass у кожному lab*/main.tex.
%
% ЗБІРКА: LuaLaTeX (не pdfLaTeX). Причина: пакет listings під pdfLaTeX не читає UTF-8 і зупиняє збірку на кириличних коментарях у лістингах програмного коду. LuaLaTeX обробляє UTF-8 нативно.
%
% ШРИФТ: Liberation Serif — метрично сумісний з Times New Roman і містить повний набір кириличних гліфів (TeX Gyre Termes кирилиці не має, через що український текст не відображався взагалі).

\usepackage{fontspec}
\setmainfont{Liberation Serif}
\setmonofont{DejaVu Sans Mono}[Scale=0.82]

\usepackage{polyglossia}
\setmainlanguage{ukrainian}
\setotherlanguage{english}

\usepackage{microtype}

% Поля, інтервал, відступи й заголовки нижче взяті з шаблону НУЛП, за яким написана бакалаврська робота (../bachelor-thesis/thesis/main.tex). Той самий формат приймає В.М. Шиманський і для магістерських, тож усі звіти тримаємо в ньому, а не в тому, що переказували на консультації.
\usepackage[left=25mm,right=15mm,top=15mm,bottom=23mm,footskip=1.5mm]{geometry}
\usepackage{setspace}
% Саме \setstretch{1.5}, а не \onehalfspacing: останній дає близько 1.25.
\setstretch{1.5}

\usepackage{graphicx}
\usepackage{float}
\usepackage{booktabs}
\usepackage{array}
\usepackage{tabularx}
\usepackage{longtable}
\usepackage{enumitem}
% labelindent вирівнює маркер списку з абзацним відступом — так у шаблоні НУЛП.
\setlist{nosep}
\setlist[itemize]{label=\textbullet,leftmargin=*,labelindent=1.27cm}
\setlist[enumerate]{leftmargin=*,labelindent=1.27cm}
\usepackage{indentfirst}
\usepackage{amsmath}

\usepackage{hyperref}
\hypersetup{colorlinks=true,linkcolor=black,urlcolor=blue,citecolor=black}

% xurl дозволяє \url розривати адресу в будь-якому місці, а не лише після скісної риски. Без нього довгий сегмент шляху (напр. «scopus---podrobnaya- instruktsiya-dlya-uchenogo») виходить за межі сторінки, бо дефіс для hyperref не є точкою розриву. Вантажиться після hyperref — перевизначає його \UrlBreaks.
\usepackage{xurl}

% ---------------------------- Список використаної літератури ----------------------------
% biblatex + biber зі стилем IEEE: посилання «[1]» у тексті, записи — за порядком першої згадки. Обрано IEEE, а не gost-numeric (ДСТУ 8302:2015), — так вирішив автор звітів після порівняння обох варіантів.
% Кожен звіт із джерелами сам підключає свій references.bib через \addbibresource і друкує список \printbibliography[heading=none] під власним \section. Звіт без джерел від завантаженого biblatex нічого не втрачає.
% Оточення списку — стандартне для стилю: мітка від лівого поля, текст опису вирівняний після неї. Не підганяємо його під абзацний відступ нумерованих переліків тексту (labelindent у \setlist вище) — список літератури є окремим розділом зі своїм стандартним оформленням.
\usepackage[backend=biber,style=ieee]{biblatex}
% biblatex за замовчуванням додає розтяжний проміжок після «/ . : ? =» в адресах, і з вимкненими переносами рядок вирівнюється розрідженням самої адреси («https : / / www . scopus»). Рвати адресу й так дозволяє xurl, тож проміжок лише шкодить.
\setlength{\biburlbigskip}{0mu}
% Без додаткового інтервалу між описами: розділ має йти суцільним текстом, як решта звіту. Так само в шаблоні бакалаврської.
\setlength{\bibitemsep}{0pt}
\setlength{\bibparsep}{0pt}

% ---------------------------- Лістинги програмного коду ---------------------------- Використано minted (через Pygments), а не listings: listings не вміє коректно обробляти багатобайтові символи й «висмикує» кириличні коментарі з коду, зліплюючи їх в один рядок. Потребує прапорця -shell-escape.
\usepackage{xcolor}
\usepackage[newfloat]{minted}
% labelsep=period дає «Рис. 2.1. Назва» замість «Рис. 2: Назва» — формат НУЛП.
\usepackage[labelsep=period]{caption}

% baselinestretch=1 обов'язковий: без нього лістинг успадковує півторашку основного тексту, і довгий фрагмент коду перестає вміщатися у сторінку. Усередині [H]-флоата це не помилка збірки, а нескінченний цикл порожніх сторінок — lab1 так намолотив 161 тисячу сторінок.
\setminted{
  baselinestretch=1,
  fontsize=\footnotesize,
  linenos=true,
  frame=single,
  framesep=4pt,
  breaklines=true,
  breakanywhere=true,
  tabsize=4,
  numbersep=7pt,
  bgcolor=gray!8,
}
\setminted[python]{python3=true}

\SetupFloatingEnvironment{listing}{name=Лістинг,listname=Перелік лістингів}

% ---------------------------- Нумерація та оформлення заголовків ----------------------------
\usepackage{titlesec}

\setcounter{secnumdepth}{3}
\renewcommand{\thesection}{\arabic{section}}
\renewcommand{\thesubsection}{\arabic{section}.\arabic{subsection}}
\renewcommand{\thesubsubsection}{\arabic{section}.\arabic{subsection}.\arabic{subsubsection}}

% Розділ по центру, 16 пт, великими літерами; слово «РОЗДІЛ» опущене, бо звіт не дисертація. Номер лишається: \numberwithin прив'язує до нього нумерацію рисунків («Рис. 2.1»), і без нумерованих розділів вона б зламалася. \MakeUppercase четвертим аргументом отримує текст заголовка, тому в джерелі заголовки пишуться звичайним регістром.
\titleformat{\section}[block]
  {\centering\bfseries\fontsize{16}{19.2}\selectfont}
  {\thesection.}{0.3em}{\MakeUppercase}
% Підрозділи — зліва з абзацного відступу, звичайним регістром.
\titleformat{\subsection}[block]{\bfseries}{\thesubsection.}{1em}{}
\titleformat{\subsubsection}[block]{\bfseries\itshape}{\thesubsubsection.}{1em}{}

% Незіркові \titlespacing — обов'язкові. Без них titlesec застосовує зірковий варіант, який прибирає абзацний відступ у першого абзацу після заголовка й перекриває indentfirst. Значення — типові для report.
\titlespacing{\section}{0pt}{20pt}{6pt}
\titlespacing{\subsection}{1.27cm}{6pt}{0pt}
\titlespacing{\subsubsection}{1.27cm}{6pt}{0pt}

\setlength{\parindent}{1.27cm}
\setlength{\parskip}{0pt}
\renewcommand{\arraystretch}{1.3}

% Переноси слів вимкнені — так у шаблоні НУЛП. Натомість абзац розганяється міжслівними пробілами, тому tolerance і emergencystretch підняті: без них рядки виходили б за праве поле.
\hyphenpenalty=10000
\exhyphenpenalty=10000
\tolerance=9999
\setlength{\emergencystretch}{3em}

% ---------------------------- Підписи до рисунків і таблиць ----------------------------
\captionsetup[figure]{name=Рис.,justification=centering,font={bf,it}}
\captionsetup[table]{name=Таблиця,justification=raggedleft,singlelinecheck=false,labelfont=bf,textfont=normalfont}

% Рисунки, таблиці й формули нумеруються в межах розділу: «Рис. 2.1», а не «Рис. 5». Так у шаблоні НУЛП, і так само виглядали приклади на консультації.
\numberwithin{equation}{section}
\numberwithin{figure}{section}
\numberwithin{table}{section}

% ---------------------------- Колонтитули ----------------------------
\usepackage{fancyhdr}
\pagestyle{fancy}
\fancyhf{}
\fancyfoot[R]{\thepage}
\renewcommand{\headrulewidth}{0pt}
\fancypagestyle{plain}{%
  \fancyhf{}
  \fancyfoot[R]{\thepage}
  \renewcommand{\headrulewidth}{0pt}
}

\input{../../../../titlepage}
: \renewcommand{\WorkKind}{ПРАКТИЧНА РОБОТА}.
\newcommand{\WorkKind}{ЛАБОРАТОРНА РОБОТА}

\newcommand{\MakeLabTitle}[2]{%
\begin{titlepage}
% Титулка — не суцільний текст, а фіксована верстка з явних \\[...] проміжків,
% тож міжрядковий інтервал основного тексту (\setstretch{1.5} у preamble.tex)
% до неї не застосовний: він розтягував кожен із ~15 рядків на 20 %, і блок
% «Виконав» разом із «Львів --- рік» переповнював сторінку. Інтервал локальний
% для оточення titlepage, тому на решту документа це не впливає.
\setstretch{1}
\begin{center}
  {\large\bfseries МІНІСТЕРСТВО ОСВІТИ І НАУКИ УКРАЇНИ}\\[6pt]
  {\large\bfseries Національний університет «Львівська політехніка»}\\[6pt]
  {\large\bfseries Інститут комп'ютерних наук та інформаційних технологій}\\[6pt]
  {\large\bfseries Кафедра «Системи штучного інтелекту»}\\[1.2cm]

  \includegraphics[width=5cm]{../../../../logo}\\[1.2cm]

  {\Large\bfseries \WorkKind{} №#1}\\[0.5cm]
  з дисципліни\\[4pt]
  «\SubjectName»\\[0.6cm]
  на тему\\[4pt]
  \textit{#2}
\end{center}

% Відступ до блоку «Виконав» гнучкий, а не фіксований (було 2 cm). Фіксований
% разом із довгим родом роботи чи темою на два рядки переповнював сторінку, і
% «Львів --- рік» зі свого \vfill з'їжджав на окрему другу сторінку. \vfill
% стискається до нуля, коли місця бракує, тож переповнення стало неможливим.
\vfill

\begin{flushright}
  {\bfseries Виконав:}\\
  студент групи \StudentGroup\\
  \StudentName\\[1cm]
  {\bfseries Викладач:}\\
  \LecturerName\\
\end{flushright}

\vfill
\begin{center}
  Львів --- \AcademicYear
\end{center}
\end{titlepage}
\clearpage
}


% Subject metadata, read by the shared title page (titlepage.tex).

\newcommand{\SubjectName}{Методологія та процеси моделювання у штучному інтелекті}
% The lecturer's position and degree are not known yet; the other subjects put them on a line above the name, separated by `\\`.
\newcommand{\LecturerName}{Косаревич Р.Я.}
\newcommand{\StudentName}{Галич Максим}
\newcommand{\StudentGroup}{КНСШ-13}
\newcommand{\AcademicYear}{2026}

% Figures are generated artifacts, so they live in .cache/ rather than beside the text (see ARTIFACTS in src/mk1/config.py).
\graphicspath{{../../../../.cache/reports-artifacts/1th term 2026 Autumn/Методологія та процеси моделювання у штучному інтелекті/}}
% A centred column of fixed width: the wide tables below need their headers to wrap, which a plain c column cannot do.
\newcolumntype{C}[1]{>{\centering\arraybackslash}p{#1}}

\begin{document}

\MakeLabTitle{4}{Побудова та навчання нейронної мережі}

\textbf{Мета роботи.} Оволодіти методами проєктування спеціалізованих нейромережевих архітектур (зокрема автоенкодерів та екстракторів ознак) для задач детекції неоднорідностей і аномалій у режимі unsupervised / self-supervised, а також дослідити компроміс між точністю та обчислювальною складністю моделей.

Усі обчислення виконано мовою Python із бібліотеками PyTorch [\ref{ref:pytorch}], timm [\ref{ref:timm}] та Anomalib [\ref{ref:anomalib}]. Код продовжує пакет \texttt{mk1} з лабораторних робіт №1--3. Завантаження й розбиття даних описано в модулі \texttt{mvtec.py}, вилучення ознак --- у \texttt{patches.py}, запропоновану модель та її навчання --- у \texttt{feature\_ae.py}, дослідження гіперпараметрів --- у \texttt{ae\_search.py}, метрики якості --- у \texttt{anomaly\_metrics.py}, вимірювання обчислювальної складності --- у \texttt{profiling.py}, порівняння з SOTA-моделями --- у \texttt{benchmark.py}. Сценарій \texttt{lab4/run.py} лише запускає три етапи (дані, пошук гіперпараметрів, порівняння) і перевіряє, що кожен записав свої рисунки та метрики. Увесь експеримент відтворюється однією командою з фіксованим зерном генератора випадкових чисел.

\section{Дані та їх підготовка}
\label{sec:data}

\subsection{Набір даних}

Використано набір \textbf{MVTec AD} [\ref{ref:mvtec}] --- еталонний набір для детекції аномалій на промислових зображеннях без учителя. Його навчальна частина містить лише зображення без дефектів, а тестова --- зображення без дефектів і з дефектами, причому для кожного дефекту є попіксельна маска. Отже, модель навчається тільки на нормі й має знайти те, чого ніколи не бачила, а маски дають змогу оцінити не лише, \emph{чи} є дефект, а й \emph{де} він.

Обрано категорію \textbf{bottle} (пляшка, вигляд згори): найпростішу категорію набору, з одним об'єктом у центрі кадру на однорідному тлі. Вона має три типи дефектів: великий скол, малий скол і забруднення (рис.~\ref{fig:samples}). Зображення мають розмір $900 \times 900$ пікселів, формат RGB.

\begin{figure}[H]
  \centering
  \includegraphics[width=\linewidth]{lab4/samples}
  \caption{Приклад зображення кожного типу з валідаційної або тестової частини. Унизу червоним обведено маску дефекту і вказано її площу}
  \label{fig:samples}
\end{figure}

Завдання вимагає складного набору з високою дисперсією тла або незбалансованістю цілей. Категорія bottle має однорідне тло, але сильно незбалансовані цілі на рівні пікселів (розділ~\ref{sec:imbalance}): дефект займає в середньому кілька відсотків зображення, а тому аномальні пікселі становлять меншість навіть серед зображень із дефектами.

Архів набору не завантажується кодом. Офіційний сайт MVTec видає посилання лише після реєстраційної форми, а під час виконання роботи він був недоступний через помилку TLS-з'єднання. Тому категорію bottle взято з дзеркала на Hugging Face (\texttt{foersben/mvtec-ad}) під тією самою ліцензією CC BY-NC-SA 4.0. Кількість файлів у кожній папці збігається з офіційною: 209 навчальних зображень, 83 тестових і 63 маски.

\subsection{Розбиття на вибірки}

Набір MVTec AD не має валідаційної частини. Якщо обирати гіперпараметри за тестовими зображеннями, підсумкова оцінка буде завищена: у лабораторній роботі №3 показано, що найкраща оцінка пошуку містить оптимізм вибору. Тому тестову частину набору поділено навпіл зі стратифікацією за типом дефекту (лістинг~\ref{lst:split}). Перша половина --- \textbf{валідаційна}: за нею обрано всі гіперпараметри, найкращу епоху навчання та пороги для F1-Score. Друга половина --- \textbf{тестова}: її використано один раз, в остаточному порівнянні моделей (розділ~\ref{sec:benchmark}). Розподіл зображень наведено в табл.~\ref{tab:split}.

\begin{table}[H]
\caption{Кількість зображень у вибірках}
\label{tab:split}
\begin{tabularx}{\linewidth}{|X|C{2.3cm}|C{2.3cm}|C{2.3cm}|C{2.9cm}|C{1.8cm}|}
\hline
\textbf{Вибірка} & \textbf{Без дефекту} & \textbf{Великий скол} & \textbf{Малий скол} & \textbf{Забруднення} & \textbf{Усього} \\
\hline
Навчальна & 209 & --- & --- & --- & 209 \\
\hline
Валідаційна & 10 & 10 & 11 & 10 & 41 \\
\hline
Тестова & 10 & 10 & 11 & 11 & 42 \\
\hline
\end{tabularx}
\end{table}

\inputminted[firstline=179,lastline=204]{python}{../src/mk1/mvtec.py}
\captionof{listing}{Розбиття тестової частини і завантаження зображення (\texttt{src/mk1/mvtec.py})}
\label{lst:split}

\subsection{Перетворення зображень}

Кожне зображення зменшено з $900 \times 900$ до $256 \times 256$ пікселів білінійною інтерполяцією і нормалізовано середнім та стандартним відхиленням ImageNet за кожним каналом кольору. Цей розмір і цю нормалізацію використовують статті PaDiM і PatchCore: на зображеннях такого масштабу навчено переднавчені мережі. Маски зменшено до того самого розміру методом найближчого сусіда, бо інтерполяція створила б частково аномальні пікселі, яких у розмітці немає. Однакові перетворення застосовано до всіх моделей, тож порівняння в розділі~\ref{sec:benchmark} не залежить від підготовки даних. Аугментацію не застосовано: модель має вивчити саме норму, а повороти чи зміни кольору могли б зробити нормою те, що в цьому наборі є дефектом.

\subsection{Незбалансованість цілей}
\label{sec:imbalance}

Площу кожного дефекту обчислено як частку пікселів маски від площі зображення (рис.~\ref{fig:area}). Медіана становить 5,7~\%, а діапазон --- від 0,57~\% до 27,4~\%. Малі сколи найменші (медіана 2,4~\%, від 0,9~\% до 8,1~\%), великі сколи --- найбільші (медіана 11,9~\%). Забруднення мають найширший діапазон: від 0,6~\% до 16,5~\%. Загалом аномальні пікселі становлять 5,95~\% валідаційної і 5,61~\% тестової вибірки. Отже, на рівні пікселів співвідношення класів близько 1~:~16, і метрика, яка не враховує цю незбалансованість, переоцінює модель (розділ~\ref{sec:metrics}).

\begin{figure}[H]
  \centering
  \includegraphics[width=\linewidth]{lab4/defect_area}
  \caption{Площа дефекту як частка площі зображення для кожного з 63 дефектних зображень валідаційної і тестової вибірок. Вертикальна риска --- медіана типу дефекту}
  \label{fig:area}
\end{figure}

\subsection{Розбиття на патчі та вилучення ознак}
\label{sec:patches}

Ознаки вилучено переднавченими на ImageNet згортковими мережами, ваги яких не змінюються під час навчання. Карта ознак мережі з кроком 8 ділить зображення $256 \times 256$ на сітку $32 \times 32$ патчі: кожна комірка описує ділянку $8 \times 8$ пікселів (рис.~\ref{fig:patches}). Як у PatchCore [\ref{ref:patchcore}], з кожної мережі взято два шари: шар із кроком 8 зберігає локальну текстуру, а шар із кроком 16 додає контекст. Його карту збільшено білінійною інтерполяцією до сітки $32 \times 32$ і приєднано до першої за каналами. Порівняно три мережі різної складності (табл.~\ref{tab:backbones}).

\begin{table}[H]
\caption{Мережі-екстрактори ознак}
\label{tab:backbones}
\small
\begin{tabularx}{\linewidth}{|X|C{3.7cm}|C{2.9cm}|C{2.4cm}|C{1.8cm}|}
\hline
\textbf{Мережа} & \textbf{Шари (крок 8 і 16)} & \textbf{Каналів} & \textbf{Параметрів, млн} & \textbf{GFLOPs} \\
\hline
ResNet-18 [\ref{ref:resnet}] & \texttt{layer2}, \texttt{layer3} & 128 + 256 = 384 & 2,78 & 3,66 \\
\hline
WideResNet-50-2 [\ref{ref:wrn}] & \texttt{layer2}, \texttt{layer3} & 512 + 1024 = 1536 & 24,86 & 24,00 \\
\hline
MobileNetV3-Large [\ref{ref:mobilenet}] & \texttt{blocks.2}, \texttt{blocks.4} & 40 + 112 = 152 & 2,97 & 0,56 \\
\hline
\end{tabularx}
\end{table}

Кількість параметрів і FLOPs виміряно на одне зображення $256 \times 256$ для мережі в тому вигляді, у якому її створює бібліотека. Ваги завантажено з репозиторію timm [\ref{ref:timm}]. ResNet-18 і WideResNet-50-2 timm обрізає після останнього потрібного шару, а MobileNetV3 --- ні: вона обчислює всі сім блоків, хоча використовуються виходи лише блоків 2 і 4. Потрібна частина має 0,79 млн параметрів і 0,40 GFLOPs, тож 73~\% параметрів і 29~\% FLOPs MobileNetV3 витрачаються даремно. Далі наведено виміряні значення, тобто вартість MobileNetV3 у цій роботі завищена.

\begin{figure}[H]
  \centering
  \includegraphics[width=0.62\linewidth]{lab4/patches}
  \caption{Сітка патчів, яку задає карта ознак із кроком 8: $32 \times 32$ комірки по $8 \times 8$ пікселів. Кольором показано окіл, за яким усереднюється вектор ознак за розміру патча $p = 1$, 3 і 5}
  \label{fig:patches}
\end{figure}

\textbf{Розмір патча} $p$ задає окіл комірок, за яким усереднюється кожен шар перед об'єднанням (лістинг~\ref{lst:aggregate}). За $p = 3$ вектор описує ділянку $24 \times 24$ пікселі, за $p = 5$ --- $40 \times 40$ пікселів, без урахування рецептивного поля мережі. Більший окіл стійкіший до шуму і незначних зсувів об'єкта, але розмиває малі дефекти. Під час усереднення крайні комірки враховують лише наявних сусідів (\texttt{count\_include\_pad=False}), щоб нулі доповнення не зменшували ознаки на краю зображення.

\inputminted[firstline=98,lastline=130]{python}{../src/mk1/patches.py}
\captionof{listing}{Усереднення за окілом патча і об'єднання двох шарів (\texttt{src/mk1/patches.py})}
\label{lst:aggregate}

Мережа --- найдорожча частина обчислень, і вона не змінюється. Тому під час навчання її виходи для всіх зображень обчислено один раз, а усереднення за окілом $p$ застосовується до збережених ознак на кожному кроці. Так одне проходження мережі обслуговує всі розміри патча та всі епохи.

\section{Архітектура запропонованої моделі}
\label{sec:model}

\subsection{Ідея: автоенкодер ознак}

Запропоновано \textbf{автоенкодер, що відновлює ознаки патчів}, а не пікселі, у підході DFR (Deep Feature Reconstruction) [\ref{ref:dfr}]. Автоенкодер навчається пропускати вектор ознак кожного патча крізь вузьке горло й відновлювати його. Відновлювати він уміє лише те, що бачив, тобто патчі без дефектів. Патч із дефектом відновлюється погано, і похибка відновлення є оцінкою аномальності.

Відновлення ознак замість пікселів --- принципове рішення. Піксельний автоенкодер розмиває дрібні деталі, тож похибка відновлення висока і на нормальній текстурі: модель «бачить» аномалії всюди, де є дрібний візерунок. Автоенкодер ознак працює в просторі, де мережа вже відкинула шум та освітлення і зберегла опис вмісту.

\subsection{Архітектура}

Автоенкодер повністю згортковий (лістинг~\ref{lst:ae}, табл.~\ref{tab:ae}). Згортки $1 \times 1$ на краях стискають кожен вектор окремо, а згортка $3 \times 3$ з кожного боку горла відновлює патч з урахуванням сусідів. Отже, дефект не можна відтворити, просто скопіювавши патч через горло: відновлення спирається на окіл, який здебільшого нормальний. Горло має 64 канали, тобто в 2,4 раза вужче за вхід навіть для MobileNetV3 і в 24 рази --- для WideResNet-50-2. Після кожної прихованої згортки йдуть пакетна нормалізація [\ref{ref:bn}] і ReLU.

\begin{table}[H]
\caption{Шари автоенкодера ($C$ --- кількість каналів ознак, сітка $32 \times 32$)}
\label{tab:ae}
\small
\begin{tabularx}{\linewidth}{|c|X|c|c|}
\hline
\textbf{№} & \textbf{Шар} & \textbf{Вхід $\rightarrow$ вихід} & \textbf{Параметрів за $C = 152$} \\
\hline
1 & згортка $1 \times 1$, BatchNorm, ReLU & $C \rightarrow 256$ & 39\,424 \\
\hline
2 & згортка $3 \times 3$, BatchNorm, ReLU & $256 \rightarrow 128$ & 295\,168 \\
\hline
3 & згортка $1 \times 1$ (горло) & $128 \rightarrow 64$ & 8\,256 \\
\hline
4 & згортка $1 \times 1$, BatchNorm, ReLU & $64 \rightarrow 128$ & 8\,448 \\
\hline
5 & згортка $3 \times 3$, BatchNorm, ReLU & $128 \rightarrow 256$ & 295\,424 \\
\hline
6 & згортка $1 \times 1$ & $256 \rightarrow C$ & 39\,064 \\
\hline
\multicolumn{3}{|r|}{\textbf{Усього}} & \textbf{685\,784} \\
\hline
\end{tabularx}
\end{table}

\inputminted[firstline=87,lastline=104]{python}{../src/mk1/feature_ae.py}
\captionof{listing}{Шари автоенкодера (\texttt{src/mk1/feature\_ae.py})}
\label{lst:ae}

Вхідні ознаки стандартизовано за кожним каналом середнім і стандартним відхиленням навчальних патчів. Канали ознак відрізняються за масштабом на порядки, і без стандартизації функцію втрат визначали б кілька найбільших. Статистики обчислено лише на навчальних зображеннях, як і будь-який крок, що навчається на даних (лабораторна робота №3).

\subsection{Оцінка аномальності та регуляризація}

Оцінка патча --- середньоквадратична похибка відновлення за всіма каналами (лістинг~\ref{lst:forward}). Карта $32 \times 32$ збільшується до $256 \times 256$ і згладжується гаусовим фільтром із $\sigma = 4$ тим самим кодом, що й у PatchCore з бібліотеки Anomalib. Тому карти двох моделей відрізняються лише способом оцінювання патча. Оцінка зображення --- максимум його карти.

Щоб автоенкодер не узагальнював настільки добре, що відновлює й дефекти, використано два регуляризатори:

\begin{itemize}
\item \textbf{weight decay} --- розщеплене згасання ваг оптимізатора AdamW [\ref{ref:adamw}];
\item \textbf{маскувальний шум} --- під час навчання частка вхідних ознак обнуляється, і автоенкодер має відновити їх з решти (шумопонижувальний автоенкодер [\ref{ref:dae}]). Значення не масштабуються, як у dropout: обнулені ознаки --- це пошкодження, яке треба виправити, а не випадковість, яку слід усереднити.
\end{itemize}

\inputminted[firstline=110,lastline=138]{python}{../src/mk1/feature_ae.py}
\captionof{listing}{Відновлення з маскувальним шумом і оцінка патча (\texttt{src/mk1/feature\_ae.py})}
\label{lst:forward}

\subsection{Навчання}

Автоенкодер навчається оптимізатором AdamW із кроком $10^{-3}$, функція втрат --- середньоквадратична похибка відновлення стандартизованих ознак, 30 епох, пакет із 8 зображень (тобто 8\,192 векторів патчів). Після кожної епохи модель оцінюється на валідаційній вибірці двома показниками (лістинг~\ref{lst:fit}):

\begin{itemize}
\item \textbf{валідаційні втрати} --- похибка відновлення на 10 валідаційних зображеннях без дефектів. Вона показує, чи узагальнює автоенкодер норму на зображення, яких не бачив;
\item \textbf{Pixel\_AUROC} на всій валідаційній вибірці --- чи відрізняє автоенкодер дефекти від норми.
\end{itemize}

Повертається модель тієї епохи, на якій валідаційний Pixel\_AUROC найвищий, тобто рання зупинка, обрана лише за валідаційною вибіркою.

\inputminted[firstline=287,lastline=316]{python}{../src/mk1/feature_ae.py}
\captionof{listing}{Цикл навчання з оцінкою після кожної епохи (\texttt{src/mk1/feature\_ae.py})}
\label{lst:fit}

\section{Метрики якості та обчислювальної складності}
\label{sec:metrics}

Кожна модель --- і запропонована, і SOTA-моделі з Anomalib --- повертає одне й те саме: карту оцінок аномальності, по одній оцінці на піксель. Усі карти оцінюються тим самим кодом (модуль \texttt{anomaly\_metrics.py}), тож порівняння не залежить від реалізації метрик.

\begin{itemize}
\item \textbf{Pixel\_AUROC} --- площа під ROC-кривою за всіма пікселями вибірки відносно масок. Це метрика, яку вимагає завдання і наводять статті про MVTec AD.
\item \textbf{Pixel PR-AUC} (average precision) --- те саме ранжування, оцінене за кривою точність-повнота. Аномальні пікселі становлять 5--6~\%, і ROC-AUC мало чутлива до того, скільки нормальних пікселів випереджають аномальні: на кожен аномальний припадає 16 нормальних, і навіть мала частка хибних тривог серед них --- це багато хибних пікселів. PR-AUC, основна метрика лабораторних робіт №1--3, це враховує [\ref{ref:saito}].
\item \textbf{F1-Score} на рівні пікселів і зображень. F1 потребує порогу, а поріг --- це вибір. Його встановлено там, де F1 найвищий на \emph{валідаційній} вибірці, і без змін застосовано до тестової (лістинг~\ref{lst:score}). Тестовий F1 --- це якість детектора, який тестових зображень не бачив.
\item \textbf{Image\_AUROC} --- ранжування зображень за максимумом їхньої карти.
\end{itemize}

\inputminted[firstline=105,lastline=120]{python}{../src/mk1/anomaly_metrics.py}
\captionof{listing}{Пороги F1 з валідаційної вибірки, застосовані до тестової (\texttt{src/mk1/anomaly\_metrics.py})}
\label{lst:score}

Обчислювальну складність виміряно для детектора повністю, від нормалізованого зображення до карти аномальності, бо саме це оплачує застосування: у PatchCore більшу частину часу займає пошук найближчих сусідів після мережі.

\begin{itemize}
\item \textbf{FLOPs} пораховано засобом \texttt{FlopCounterMode} бібліотеки PyTorch, що реєструє кожне множення матриць і кожну згортку, тож пошук сусідів і відстань Махаланобіса враховано разом із мережами. Одне множення з додаванням --- дві операції. Поелементні операції не враховано, вони становлять малу частку в усіх моделях.
\item \textbf{Затримка} --- реальний час одного проходження на CPU з пакетом з одного зображення, як вони надходять на лінії контролю якості: 5 прогрівальних запусків, потім медіана, 10-й і 90-й процентилі за 30 запусками (лістинг~\ref{lst:latency}). Кількість потоків зафіксовано на 8, по одному на фізичне ядро процесора AMD Ryzen 7 8845HS.
\end{itemize}

GPU, придатного для PyTorch, на комп'ютері немає: вбудовану графіку Radeon 780M використана збірка PyTorch не підтримує. Тому всі вимірювання виконано на CPU. Для кожного детектора окремо виміряно й саму мережу-екстрактор: різниця з повним детектором --- вартість того, що працює після мережі.

\inputminted[firstline=103,lastline=121]{python}{../src/mk1/profiling.py}
\captionof{listing}{Вимірювання затримки на CPU (\texttt{src/mk1/profiling.py})}
\label{lst:latency}

\section{Дослідження гіперпараметрів}
\label{sec:search}

\subsection{Схема пошуку}

Досліджено чотири гіперпараметри: мережу-екстрактор, розмір патча $p$ і силу обох регуляризаторів. Повний перебір тих самих значень (табл.~\ref{tab:space}) потребував би $3 \cdot 3 \cdot 4 \cdot 3 = 108$ навчань. Тому пошук виконано \textbf{по одному фактору}: кожен крок змінює один гіперпараметр, починаючи з переможця попереднього кроку. Такий пошук потребує лише 10 навчань, бо переможець попереднього кроку повторно не навчається. Він припускає, що фактори взаємодіють слабко, тому першим іде фактор, який змінює самі ознаки, --- мережа-екстрактор. Початкова конфігурація: ResNet-18, $p = 3$, weight decay $10^{-2}$ (типове значення AdamW), без маскувального шуму.

\begin{table}[H]
\caption{Значення гіперпараметрів у пошуку}
\label{tab:space}
\begin{tabularx}{\linewidth}{|c|X|X|}
\hline
\textbf{Крок} & \textbf{Гіперпараметр} & \textbf{Значення} \\
\hline
1 & мережа-екстрактор & ResNet-18, WideResNet-50-2, MobileNetV3-Large \\
\hline
2 & розмір патча $p$ & 1, 3, 5 \\
\hline
3 & weight decay & 0; 0,01; 0,1; 1 \\
\hline
4 & частка маскованих ознак & 0; 0,2; 0,5 \\
\hline
\end{tabularx}
\end{table}

Кожну конфігурацію оцінено за найкращою епохою на валідаційній вибірці. Тестова вибірка на цьому етапі взагалі не завантажується.

\subsection{Результати}

Результати наведено в табл.~\ref{tab:search} та на рис.~\ref{fig:search}. Увесь пошук тривав 15 хвилин.

\begin{table}[H]
\caption{Результати пошуку на валідаційній вибірці (жирним --- переможець кроку)}
\label{tab:search}
\small
\renewcommand{\arraystretch}{1.15}
\begin{tabularx}{\linewidth}{|X|C{2.3cm}|C{2.1cm}|C{2.3cm}|C{1.7cm}|C{1.3cm}|}
\hline
\textbf{Конфігурація} & \textbf{Pixel\_\allowbreak AUROC} & \textbf{Pixel PR-\allowbreak AUC} & \textbf{Image\_\allowbreak AUROC} & \textbf{Найкр. епоха} & \textbf{Час, с} \\
\hline
\multicolumn{6}{|l|}{\textit{Крок 1: мережа ($p = 3$, weight decay 0,01, без шуму)}} \\
\hline
ResNet-18 & 0,9486 & 0,3832 & 0,9935 & 25 & 83 \\
\hline
WideResNet-50-2 & 0,9779 & 0,6256 & 0,9968 & 3 & 172 \\
\hline
\textbf{MobileNetV3-Large} & \textbf{0,9848} & \textbf{0,7521} & \textbf{0,9968} & \textbf{4} & \textbf{73} \\
\hline
\multicolumn{6}{|l|}{\textit{Крок 2: розмір патча (MobileNetV3-Large)}} \\
\hline
$p = 1$ & 0,9832 & 0,7301 & 1,0000 & 8 & 67 \\
\hline
\textbf{$p = 3$} & \textbf{0,9848} & \textbf{0,7521} & \textbf{0,9968} & \textbf{4} & \textbf{73} \\
\hline
$p = 5$ & 0,9753 & 0,6517 & 1,0000 & 3 & 68 \\
\hline
\multicolumn{6}{|l|}{\textit{Крок 3: weight decay ($p = 3$)}} \\
\hline
\textbf{0} & \textbf{0,9848} & \textbf{0,7521} & \textbf{0,9968} & \textbf{4} & \textbf{67} \\
\hline
0,01 & 0,9848 & 0,7521 & 0,9968 & 4 & 73 \\
\hline
0,1 & 0,9848 & 0,7559 & 0,9968 & 3 & 68 \\
\hline
1 & 0,9847 & 0,7561 & 0,9968 & 3 & 67 \\
\hline
\multicolumn{6}{|l|}{\textit{Крок 4: частка маскованих ознак (weight decay 0)}} \\
\hline
0 & 0,9848 & 0,7521 & 0,9968 & 4 & 67 \\
\hline
\textbf{0,2} & \textbf{0,9851} & \textbf{0,7610} & \textbf{0,9968} & \textbf{3} & \textbf{71} \\
\hline
0,5 & 0,9763 & 0,7010 & 0,9871 & 22 & 74 \\
\hline
\end{tabularx}
\end{table}

\begin{figure}[H]
  \centering
  \includegraphics[width=\linewidth]{lab4/search}
  \caption{Pixel\_AUROC (угорі) і Pixel PR-AUC (унизу) кожної конфігурації на валідаційній вибірці, по стовпцю на крок пошуку. Колом позначено переможця кроку}
  \label{fig:search}
\end{figure}

\textbf{Мережа-екстрактор --- найважливіший фактор.} Різниця між мережами (0,036 Pixel\_AUROC і 0,37 PR-AUC) більша, ніж сумарний вплив решти трьох гіперпараметрів. Найкращою виявилася найменша мережа, MobileNetV3-Large, хоча вона має у 43 рази менше FLOPs, ніж WideResNet-50-2. Найгіршою --- ResNet-18. Автоенкодер на її ознаках повільно покращувався протягом усіх 30 епох (з 0,9385 до 0,9486 на 25-й епосі), але не досяг рівня інших мереж. Причина не в ознаках ResNet-18 як таких: PaDiM і PatchCore на тих самих вагах мають Pixel\_AUROC 0,978 (розділ~\ref{sec:benchmark}). Карти аномальності (рис.~\ref{fig:maps}) показують, що автоенкодер на ResNet-18 дуже сильно реагує на великі сколи, але майже не реагує на малі сколи та забруднення. Отже, на цих ознаках похибка відновлення має різний масштаб для різних типів дефектів, і за спільного ранжування пікселів малі дефекти програють великим. Найсильніше це видно за PR-AUC (0,38 проти 0,75).

\textbf{Розмір патча $p = 3$ найкращий}, а $p = 5$ --- найгірший (0,9753). За медіанної площі 2,4~\% малий скол займає близько $40 \times 40$ пікселів зображення $256 \times 256$, тобто стільки ж, скільки окіл $p = 5$. Половина малих сколів менша, і за $p = 5$ їхні ознаки усереднюються з нормальним оточенням. Окіл $p = 1$ втрачає небагато (0,9832), але досягає найкращого результату пізніше, на восьмій епосі.

\textbf{Weight decay майже не впливає на результат.} Значення 0 і 0,01 дали однакові оцінки до четвертого знака, 0,1 і 1 --- ті самі в межах 0,0001 Pixel\_AUROC. Причина в тому, що основним регуляризатором виявилася рання зупинка. Найкраща епоха --- 3--4, тобто лише 80--110 кроків оптимізатора. AdamW зменшує ваги на частку $\eta\lambda$ за крок: за $\lambda = 0{,}01$ і кроку $\eta = 10^{-3}$ за 100 кроків ваги зменшуються на 0,1~\%, і навіть за $\lambda = 1$ --- лише на 10~\%. За такий короткий час розщеплене згасання ваг не встигає вплинути. Переможцем кроку формально став $\lambda = 0$, але його перевага над $\lambda = 0{,}01$ --- 0,0000015 Pixel\_AUROC, тобто випадкова.

\textbf{Маскувальний шум із часткою 0,2 дав невелике покращення} (+0,0003 Pixel\_AUROC, +0,009 PR-AUC), а частка 0,5 --- погіршення: автоенкодер, якому доводиться відновлювати половину входу, навчається повільно (Pixel\_AUROC після першої епохи 0,895 проти 0,983), досягає найкращого результату лише на 22-й епосі і все одно поступається.

Різниці на кроках 3 і 4 (до 0,0004 Pixel\_AUROC) того самого порядку, що й коливання між сусідніми епохами одного навчання (рис.~\ref{fig:training}), тож вибір між цими конфігураціями визначає радше шум, ніж реальна перевага. Обрана конфігурація: \textbf{MobileNetV3-Large, $p = 3$, weight decay 0, частка маскованих ознак 0,2}.

\section{Навчання обраної моделі}
\label{sec:training}

Криві навчання обраної конфігурації показано на рис.~\ref{fig:training}, ключові епохи --- у табл.~\ref{tab:training}. Одна епоха разом з оцінюванням на валідаційній вибірці триває 2,4~с на CPU, усі 30 епох --- 71~с.

\begin{figure}[H]
  \centering
  \includegraphics[width=\linewidth]{lab4/training}
  \caption{Навчання обраної конфігурації. Ліворуч --- втрати відновлення на навчальній вибірці (з маскувальним шумом) і на валідаційних зображеннях без дефектів (без шуму), праворуч --- Pixel\_AUROC на валідаційній вибірці після кожної епохи}
  \label{fig:training}
\end{figure}

\begin{table}[H]
\caption{Втрати й метрика на окремих епохах}
\label{tab:training}
\begin{tabularx}{\linewidth}{|X|c|c|c|c|c|c|}
\hline
\textbf{Епоха} & \textbf{1} & \textbf{3} & \textbf{5} & \textbf{10} & \textbf{20} & \textbf{30} \\
\hline
Втрати, навчальна & 0,356 & 0,069 & 0,050 & 0,035 & 0,026 & 0,022 \\
\hline
Втрати, валідаційна & 0,160 & 0,084 & 0,069 & 0,055 & 0,044 & 0,039 \\
\hline
Pixel\_AUROC, валідаційна & 0,9831 & \textbf{0,9851} & 0,9844 & 0,9831 & 0,9819 & 0,9811 \\
\hline
\end{tabularx}
\end{table}

\textbf{Втрати зменшуються протягом усього навчання, а якість детекції --- ні.} Обидві криві втрат спадають монотонно: навчальні --- у 16 разів, валідаційні --- у 4 рази. Pixel\_AUROC досягає максимуму вже на третій епосі (0,9851) і потім повільно знижується до 0,9811 на тридцятій. Це головна особливість автоенкодерів для детекції аномалій: модель навчається не «розпізнавати норму», а відновлювати ознаки, і з кожною епохою відновлює краще \emph{все}, зокрема й дефекти, яких не бачила. Різниця між похибкою на дефекті та на нормі зменшується, і детекція погіршується, хоча за функцією втрат модель стає кращою. Отже, функція втрат не може бути критерієм вибору моделі, і рання зупинка за валідаційною метрикою детекції є необхідною частиною методу.

Починаючи з другої епохи, валідаційні втрати вищі за навчальні, і розрив між ними зростає: з 1,1 раза на другій епосі до 1,8 раза на тридцятій. Отже, автоенкодер дещо перенавчається на навчальних зображеннях, і насправді більше, ніж показують криві: навчальні втрати обчислено з маскувальним шумом, який їх завищує. На першій епосі навчальні втрати вищі за валідаційні, бо вони усереднюються протягом епохи, за яку модель швидко покращується, а валідаційні обчислюються вже після неї.

\section{Профілювання та порівняння з SOTA-моделями}
\label{sec:benchmark}

\subsection{Моделі для порівняння}

Обрану модель порівняно з двома еталонними методами детекції аномалій у реалізації бібліотеки Anomalib [\ref{ref:anomalib}] з її типовими параметрами:

\begin{itemize}
\item \textbf{PaDiM} [\ref{ref:padim}] моделює ознаки кожної комірки сітки багатовимірним гаусовим розподілом і оцінює патч відстанню Махаланобіса. Мережа ResNet-18, шари \texttt{layer1}--\texttt{layer3}, 100 випадково обраних каналів, сітка $64 \times 64$;
\item \textbf{PatchCore} [\ref{ref:patchcore}] зберігає банк пам'яті з 10~\% навчальних патчів, відібраних жадібним алгоритмом coreset, і оцінює патч відстанню до найближчого сусіда в банку. Мережа WideResNet-50-2, як у статті, і ResNet-18 --- для порівняння з автоенкодером на тій самій мережі.
\end{itemize}

PaDiM на WideResNet-50-2 не досліджено: з 550 каналами Anomalib зберігає для кожної з $64 \times 64$ комірок матрицю коваріації $550 \times 550$, тобто понад 5~ГБ, більше, ніж має вільної пам'яті комп'ютер. Крім того, обрану конфігурацію автоенкодера (з $p$ і регуляризацією з розділу~\ref{sec:search}) навчено й на двох інших мережах, щоб відокремити внесок мережі від внеску автоенкодера. Усі моделі отримують ті самі зображення, у тому самому розмірі й нормалізації, і оцінюються на тестовій вибірці тим самим кодом. Код навчання обох базових моделей наведено в лістингу~\ref{lst:baselines}.

\inputminted[firstline=130,lastline=151]{python}{../src/mk1/benchmark.py}
\captionof{listing}{Навчання PaDiM і PatchCore з бібліотеки Anomalib (\texttt{src/mk1/benchmark.py})}
\label{lst:baselines}

\subsection{Якість}

Метрики якості на тестовій вибірці наведено в табл.~\ref{tab:quality}.

\begin{table}[H]
\caption{Якість на тестовій вибірці. Пороги F1 обрано на валідаційній вибірці}
\label{tab:quality}
\small
\renewcommand{\arraystretch}{1.15}
\begin{tabularx}{\linewidth}{|X|C{2.1cm}|C{2.0cm}|C{1.5cm}|C{2.1cm}|C{1.6cm}|}
\hline
\textbf{Модель} & \textbf{Pixel\_\allowbreak AUROC} & \textbf{Pixel PR-\allowbreak AUC} & \textbf{Pixel F1} & \textbf{Image\_\allowbreak AUROC} & \textbf{Image F1} \\
\hline
\textbf{Автоенкодер (MobileNetV3-Large)} & \textbf{0,9824} & \textbf{0,7193} & \textbf{0,7103} & \textbf{1,0000} & \textbf{0,9841} \\
\hline
Автоенкодер (WideResNet-50-2) & 0,9739 & 0,5640 & 0,6397 & 0,9969 & 0,9841 \\
\hline
Автоенкодер (ResNet-18) & 0,9470 & 0,3785 & 0,4914 & 1,0000 & 0,9841 \\
\hline
PaDiM (ResNet-18) & 0,9776 & 0,6359 & 0,6614 & 0,9969 & 0,9841 \\
\hline
PatchCore (ResNet-18) & 0,9782 & 0,6916 & 0,6716 & 1,0000 & 0,9841 \\
\hline
PatchCore (WideResNet-50-2) & 0,9848 & 0,7611 & 0,7159 & 1,0000 & 0,9841 \\
\hline
\end{tabularx}
\end{table}

\textbf{Запропонована модель поступається лише PatchCore на WideResNet-50-2}, і то небагато: на 0,0024 Pixel\_AUROC, 0,042 PR-AUC і 0,006 Pixel F1. PaDiM і PatchCore на ResNet-18 вона перевершує за всіма метриками на рівні пікселів: Pixel\_AUROC на 0,004--0,005, PR-AUC на 0,03--0,08, F1 на 0,04--0,05.

\textbf{На рівні зображень задача насичена.} Image\_AUROC усіх моделей 0,997--1,000, тобто дефектні зображення майже ідеально ранжуються вище за нормальні. Image F1 у всіх однаковий (0,9841 --- одне з 32 дефектних зображень нижче порогу): поріг, обраний на 41 валідаційному зображенні, на тестовій вибірці виявився трохи зависоким. Отже, категорія bottle не розрізняє моделі на рівні зображень, і порівнювати їх слід за якістю локалізації дефектів.

\textbf{Оцінка, за якою обрано модель, оптимістична.} На валідаційній вибірці обраний автоенкодер мав Pixel\_AUROC 0,9851, на тестовій --- 0,9824 (на 0,0027 менше). У PatchCore на WideResNet-50-2, гіперпараметри якого на цих даних не обирали, різниця між двома половинами в 2,5 раза менша (0,0011), у PaDiM її немає зовсім. Імовірно, частина валідаційної оцінки автоенкодера --- це оптимізм вибору серед 10 конфігурацій і 30 епох, той самий ефект, що й у лабораторній роботі №3, хоча на 41 і 42 зображеннях різниця такого розміру може бути й випадковою. Саме тому для порівняння використано тестову вибірку, а не валідаційну.

Карти аномальності кожної моделі для одного тестового зображення кожного типу наведено на рис.~\ref{fig:maps}.

\begin{figure}[H]
  \centering
  \includegraphics[width=\linewidth]{lab4/anomaly_maps}
  \caption{Карти аномальності на тестових зображеннях. Ліворуч --- зображення з обведеною маскою дефекту. Шкала кольорів кожної моделі --- від 1-го до 99,9-го процентиля її оцінок на всій тестовій вибірці, тож кольори можна порівнювати в межах стовпця}
  \label{fig:maps}
\end{figure}

Автоенкодер на MobileNetV3 локалізує всі три типи дефектів, і його карта нормального зображення рівномірно холодна. Карти PatchCore гладші й ширші: на зображенні з великим сколом підвищена оцінка охоплює майже все кільце горлечка навколо дефекту. Автоенкодер на ResNet-18 яскраво виділяє великий скол, але малий скол і забруднення в його шкалі майже не видно --- це і є причина його низької PR-AUC.

\subsection{Обчислювальна складність}

Обчислювальну складність наведено в табл.~\ref{tab:cost} та на рис.~\ref{fig:cost}. «Параметри» --- ваги мережі-екстрактора та навчені ваги автоенкодера, «пам'ять моделі» --- те, що метод зберігає після навчання: гаусові розподіли PaDiM і банк пам'яті PatchCore.

\begin{table}[H]
\caption{Обчислювальна складність на CPU (AMD Ryzen 7 8845HS, 8 потоків, пакет з одного зображення $256 \times 256$)}
\label{tab:cost}
\footnotesize
\renewcommand{\arraystretch}{1.15}
\setlength{\tabcolsep}{3pt}
\begin{tabularx}{\linewidth}{|X|C{2.1cm}|C{2.0cm}|C{1.7cm}|C{2.4cm}|C{1.7cm}|C{1.9cm}|}
\hline
\textbf{Модель} & \textbf{Параметри мережі / голови, млн} & \textbf{Пам'ять моделі, млн чисел} & \textbf{GFLOPs} & \textbf{Затримка, мс/зобр. (p10--p90)} & \textbf{З них мережа, мс} & \textbf{Навчання, с} \\
\hline
\textbf{Автоенкодер (MobileNetV3-Large)} & \textbf{2,97 / 0,69} & --- & \textbf{2,10} & \textbf{13,8 (13,1--18,4)} & \textbf{9,3} & \textbf{71} \\
\hline
Автоенкодер (WideResNet-50-2) & 24,86 / 1,40 & --- & 26,99 & 57,0 (51,7--62,7) & 42,5 & 219 \\
\hline
Автоенкодер (ResNet-18) & 2,78 / 0,81 & --- & 5,45 & 17,1 (16,3--18,1) & 10,0 & 94 \\
\hline
PaDiM (ResNet-18) & 2,78 / --- & 41,4 & 3,89 & 39,6 (33,6--47,8) & 7,4 & 8 \\
\hline
PatchCore (ResNet-18) & 2,78 / --- & 8,2 & 20,65 & 72,7 (71,8--76,5) & 7,1 & 895 \\
\hline
PatchCore (WideResNet-50-2) & 24,86 / --- & 32,9 & 91,53 & 247,9 (238,5--264,8) & 44,2 & 1025 \\
\hline
\end{tabularx}
\end{table}

\begin{figure}[H]
  \centering
  \includegraphics[width=\linewidth]{lab4/quality_cost}
  \caption{Якість на тестовій вибірці залежно від затримки на CPU (логарифмічна шкала). Площа маркера пропорційна кількості FLOPs; кола --- автоенкодер, квадрати --- моделі з Anomalib}
  \label{fig:cost}
\end{figure}

\textbf{Запропонована модель у 18 разів швидша за PatchCore на WideResNet-50-2} (13,8 проти 247,9~мс на зображення) і потребує у 44 рази менше обчислень (2,1 проти 91,5 GFLOPs), втрачаючи 0,0024 Pixel\_AUROC. Вона також найшвидша з усіх шести моделей: у 2,9 раза швидша за PaDiM і в 5,3 раза --- за PatchCore на ResNet-18, які поступаються їй і за якістю. На рис.~\ref{fig:cost} вона одна займає лівий верхній кут: єдина модель, точніша за неї, повільніша у 18 разів.

\textbf{Вартість PatchCore визначає не мережа, а пошук сусідів.} Мережа WideResNet-50-2 виконує 24 GFLOPs за 44~мс, решта --- 67,5 GFLOPs і 204~мс --- це порівняння кожного з 1\,024 патчів з усіма 21\,401 векторами банку пам'яті розмірністю 1\,536. Тому заміна мережі на ResNet-18 зменшує FLOPs мережі в 6,6 раза, а повну затримку PatchCore --- лише в 3,4 раза: банк пам'яті має ту саму кількість векторів, лише коротших. Вартість пошуку зростає лінійно з кількістю навчальних зображень, тоді як вартість автоенкодера від неї не залежить.

\textbf{FLOPs не визначають затримку на CPU.} MobileNetV3 виконує у 6,5 раза менше операцій, ніж ResNet-18 (0,56 проти 3,66 GFLOPs), але працює не швидше: 9,3~мс проти 7--10~мс. Її глибинні (depthwise) згортки, функції активації hard-swish і блоки squeeze-and-excitation виконують мало операцій на кожен прочитаний байт, тож процесор чекає на пам'ять, а не обчислює: MobileNetV3 виконує близько 60 GFLOP/с, ResNet-18 і WideResNet-50-2 --- від 370 до 560. Так само PaDiM виконує лише 0,23 GFLOPs після мережі, але це займає 32~мс: для кожної з 4\,096 комірок сітки потрібно прочитати власну обернену матрицю коваріації $100 \times 100$, тобто 164~МБ на кожне зображення. Перевага MobileNetV3 у запропонованій моделі полягає не в швидкості мережі, а в тому, що 152 канали її ознак роблять автоенкодер дешевим (1,5 GFLOPs і близько 4,5~мс), і в тому, що на її ознаках автоенкодер найточніший.

\textbf{Навчання.} Автоенкодер навчається 71~с (30 епох), але найкраща епоха третя, тож за ранньої зупинки досить близько 10~с. PaDiM навчається найшвидше (8~с): він лише обчислює середні та коваріації. PatchCore навчається 15--17 хвилин: жадібний вибір 21\,401 вектора з 214\,016 патчів вимагає 21\,401 проходу по всіх патчах.

\textbf{Пам'ять.} Після навчання автоенкодер зберігає 0,69 млн ваг (2,7~МБ), PaDiM --- 41,4 млн чисел гаусових розподілів (165~МБ), PatchCore --- банк пам'яті з 8,2 або 32,9 млн чисел (33 або 131~МБ). Банк PatchCore і розподіли PaDiM зростають із роздільністю сітки та, у PatchCore, з кількістю навчальних зображень, а розмір автоенкодера від них не залежить.

\section{Висновки}

У лабораторній роботі для детекції дефектів на зображеннях без учителя (категорія bottle набору MVTec AD) запропоновано, досліджено й порівняно з SOTA-рішеннями автоенкодер, що відновлює ознаки патчів переднавченої мережі.

\begin{enumerate}
\item \textbf{Стан даних і виконані перетворення.} Навчальна вибірка --- 209 зображень без дефектів, валідаційна і тестова --- по половині тестової частини набору (41 і 42 зображення), розбитої зі стратифікацією за трьома типами дефектів. Зображення зменшено з $900 \times 900$ до $256 \times 256$ пікселів і нормалізовано статистиками ImageNet, маски --- зменшено методом найближчого сусіда. Цілі сильно незбалансовані: дефект займає від 0,6~\% до 27~\% зображення (медіана 5,7~\%), аномальні пікселі становлять 5,6--6,0~\% вибірки. Тому якість оцінено не лише Pixel\_AUROC, а й Pixel PR-AUC. Зображення розбито на сітку $32 \times 32$ патчі, кожен з яких описано ознаками двох шарів переднавченої мережі, усередненими за окілом $p \times p$ комірок.

\item \textbf{Архітектура.} Повністю згортковий автоенкодер із горлом у 64 канали і згортками $3 \times 3$ навколо нього відновлює стандартизовані ознаки кожного патча з урахуванням сусідів. Похибка відновлення є оцінкою аномальності. Модель має 0,69 млн навчених параметрів.

\item \textbf{Гіперпараметри.} За 10 навчань (замість 108 у повному переборі) обрано MobileNetV3-Large, $p = 3$, без weight decay, з маскуванням 20~\% вхідних ознак. Мережа-екстрактор виявилася найважливішим фактором: різниця між мережами 0,036 Pixel\_AUROC, більша за вплив решти гіперпараметрів разом. Найкращою була найменша мережа. Окіл $p = 5$ розмиває малі сколи. Weight decay не вплинув на результат, бо головним регуляризатором стала рання зупинка після 3--4 епох.

\item \textbf{Навчання.} Втрати на навчальній і валідаційній вибірках зменшуються протягом усіх 30 епох, а Pixel\_AUROC досягає максимуму на третій епосі й потім знижується: автоенкодер учиться відновлювати й дефекти. Отже, функцію втрат не можна використовувати для вибору моделі, а рання зупинка за метрикою детекції на валідаційній вибірці --- обов'язкова.

\item \textbf{Порівняння з SOTA.} На тестовій вибірці автоенкодер має Pixel\_AUROC 0,9824, PR-AUC 0,719 і F1 0,710. Він перевершує PaDiM і PatchCore на ResNet-18 і поступається лише PatchCore на WideResNet-50-2 (0,9848, 0,761 і 0,716). На рівні зображень усі моделі майже ідеальні (Image\_AUROC 0,997--1,000), тож категорія bottle розрізняє їх лише за локалізацією дефектів.

\item \textbf{Компроміс точності та складності.} Автоенкодер обробляє зображення за 13,8~мс на CPU і виконує 2,1 GFLOPs --- у 18 разів швидше й у 44 рази дешевше, ніж PatchCore на WideResNet-50-2 (247,9~мс, 91,5 GFLOPs), за втрати 0,0024 Pixel\_AUROC. 82~\% часу PatchCore займає пошук найближчих сусідів, і ця частина зростає з обсягом навчальних даних. Затримка на CPU не пропорційна FLOPs: MobileNetV3 має в 6,5 раза менше FLOPs, ніж ResNet-18, але працює не швидше, а PaDiM з найменшою кількістю FLOPs утричі повільніший за автоенкодер, бо читає 164~МБ коваріацій на кожне зображення.
\end{enumerate}

Головний висновок роботи: для детекції дефектів на CPU доцільна архітектура, що спрямовує обчислення на навчену компактну модель, а не на порівняння з усіма навчальними прикладами. Невеликий автоенкодер на ознаках легкої мережі досягає точності, близької до найкращої SOTA-моделі, за частку її вартості, і його вартість не зростає з обсягом даних. Обмеження роботи: одна, найпростіша категорія MVTec AD, валідаційна і тестова вибірки по 41--42 зображення, тож різниці менші за 0,003 Pixel\_AUROC не є надійними. Затримку виміряно лише на CPU.

\section*{Список використаних джерел}

\begin{enumerate}[label={[\arabic*]},ref={\arabic*},leftmargin=*,labelindent=0pt]
\item \label{ref:mvtec} Bergmann P., Fauser M., Sattlegger D., Steger C. MVTec AD --- A Comprehensive Real-World Dataset for Unsupervised Anomaly Detection // Proc. IEEE/CVF Conference on Computer Vision and Pattern Recognition (CVPR). 2019. P. 9592--9600.
\item \label{ref:patchcore} Roth K., Pemula L., Zepeda J., Schölkopf B., Brox T., Gehler P. Towards Total Recall in Industrial Anomaly Detection // Proc. IEEE/CVF Conference on Computer Vision and Pattern Recognition (CVPR). 2022. P. 14318--14328.
\item \label{ref:padim} Defard T., Setkov A., Loesch A., Audigier R. PaDiM: A Patch Distribution Modeling Framework for Anomaly Detection and Localization // Pattern Recognition. ICPR International Workshops and Challenges. Lecture Notes in Computer Science. 2021. Vol. 12664. P. 475--489.
\item \label{ref:dfr} Shi Y., Yang J., Qi Z. Unsupervised Anomaly Segmentation via Deep Feature Reconstruction // Neurocomputing. 2021. Vol. 424. P. 9--22.
\item \label{ref:anomalib} Akcay S., Ameln D., Vaidya A., Lakshmanan B., Ahuja N., Genc U. Anomalib: A Deep Learning Library for Anomaly Detection // Proc. IEEE International Conference on Image Processing (ICIP). 2022. P. 1706--1710.
\item \label{ref:dae} Vincent P., Larochelle H., Bengio Y., Manzagol P.-A. Extracting and Composing Robust Features with Denoising Autoencoders // Proc. 25th International Conference on Machine Learning (ICML). 2008. P. 1096--1103.
\item \label{ref:resnet} He K., Zhang X., Ren S., Sun J. Deep Residual Learning for Image Recognition // Proc. IEEE Conference on Computer Vision and Pattern Recognition (CVPR). 2016. P. 770--778.
\item \label{ref:wrn} Zagoruyko S., Komodakis N. Wide Residual Networks // Proc. British Machine Vision Conference (BMVC). 2016. P. 87.1--87.12.
\item \label{ref:mobilenet} Howard A., Sandler M., Chu G. et al. Searching for MobileNetV3 // Proc. IEEE/CVF International Conference on Computer Vision (ICCV). 2019. P. 1314--1324.
\item \label{ref:bn} Ioffe S., Szegedy C. Batch Normalization: Accelerating Deep Network Training by Reducing Internal Covariate Shift // Proc. 32nd International Conference on Machine Learning (ICML). 2015. P. 448--456.
\item \label{ref:adamw} Loshchilov I., Hutter F. Decoupled Weight Decay Regularization // International Conference on Learning Representations (ICLR). 2019.
\item \label{ref:saito} Saito T., Rehmsmeier M. The Precision-Recall Plot Is More Informative than the ROC Plot When Evaluating Binary Classifiers on Imbalanced Datasets // PLoS ONE. 2015. Vol. 10, No. 3. e0118432.
\item \label{ref:pytorch} Paszke A., Gross S., Massa F. et al. PyTorch: An Imperative Style, High-Performance Deep Learning Library // Advances in Neural Information Processing Systems 32. 2019. P. 8024--8035.
\item \label{ref:timm} Wightman R. PyTorch Image Models [Електронний ресурс]. 2019. URL: \url{https://github.com/huggingface/pytorch-image-models}.
\end{enumerate}

\end{document}
